% ============================================================
% SECTION 4: BELOW THE BOUNDARY. Written clean 2026-07-26 (evening).
% Weight: MEDIUM-HEAVY (PAPER_STRUCTURE.md). Delivers the "no benefit of any
% kind" half of the one-sentence claim.
%
% DECISIONS TAKEN WITH MOHAMMAD BEFORE DRAFTING:
%   (1) counts, not per-cell values: the claim is delimitation
%   (2) the com-Amazon exception gets its own subsection, named as an exception
%   (3) the section opens with the seven-sparsifier rationale and its exclusions
%
% CROSS-REFERENCE WARNING: sections/05_boundary.tex is currently commented out of
% main.tex, so \ref{sec:boundary} would be undefined. This section therefore
% refers to the fixed-k result in words. WHEN SECTION 5 RETURNS, replace those
% phrasings with real references (marked below with "% XREF").
%
% Evidence, all recomputed from CSVs this session, not from SUMMARY files:
%   exp_AB/results.csv     102 matched Leiden cells, 7 arms, 7 networks; 0/102
%                          positive; per-arm best cells; fragmentation columns;
%                          28 Metis rows, 8 positive at d_avg 28.5-32.6
%   exp_L/results.csv      L-Spar 7 networks, 11 ok cells, -0.0142 .. -0.4035
%   exp_P/results.csv      63 cells, 3 algorithms, 60/63 negative vs best-of-N
%   exp_P/bestof.csv       per-network best-of-N baselines
%   exp_V/results.csv      four deployments of the verified signal
% STYLE: no em-dashes, no AI-characteristic phrasing (see STYLE.md).
% ============================================================

\subsection{What was tested, and why these methods}\label{sec:negative:arms}

Seven sparsifiers were run through the protocol of Section~\ref{sec:protocol}, chosen
so that a negative result cannot be attributed to the selection.

Two are the methods whose quality claims motivate the pipeline: local similarity
sparsification~\cite{satuluri2011local} and the degree-based
sampler~\cite{liu2023dspar}, one from each family described in
Section~\ref{sec:background}. Three more are the methods that the largest
independent benchmark of sparsification ranks highest at preserving clustering
structure, namely K-Neighbor, Local Degree and Local
Similarity~\cite{chen2024demystifying, hamann2016structure}, so that the arms
include the strongest candidates chosen by someone else. The remaining two are a
connectivity-preserving backbone, a spanning structure topped up to the retention
target either at random or by similarity, which cannot disconnect a graph by
construction and is included as an instrument rather than as a competitor. Uniform
random deletion supplies the no-information baseline wherever a selection rule has
to be shown to be doing work.

Three classes are excluded, and the exclusions bound what follows. Node-removing
methods such as $k$-core peeling are not tested, because they score their
partition on a smaller and denser vertex set, which is a different comparison
requiring the coverage control of
Section~\ref{sec:protocol:retention}. Learned sparsifiers~\cite{wu2020gsgan} are
not tested, because matching them on realized retention requires retraining per
operating point. And of the detection algorithms, the two for which the founding
accuracy claim was made, Metis and Graclus, together with MLR-MCL, have no
implementation in the library used here, so the real-network arms of this section
use four other algorithms. That last exclusion is the most consequential one in
the paper and we return to it in Section~\ref{sec:discussion}.

Before any quality measurement, one structural fact constrains the design. Four of
the seven methods cannot reach aggressive retention on a sparse graph, because
each vertex keeps at least one incident edge. Table~\ref{tab:floors} gives the
floors. On the three sparsest networks no method in this study can retain less
than about a third of the edges, so the aggressive regime in which sparsification
would pay for itself is not merely unhelpful there; it is unreachable. All arms
are therefore compared at a common operating point per network, and the third
point is set at the largest floor rather than at a round number.

\begin{table}[t]
\centering
\caption{Hard retention floors by method, as a fraction of edges. A blank marks a
method with no floor. Networks are ordered by average degree; the floor is
governed by it. Below these values the method cannot be run at all, whatever the
target.}
\label{tab:floors}
\small
\setlength{\tabcolsep}{4pt}
\begin{tabular}{@{}lrrrrr@{}}
\toprule
Network & $d_{\mathrm{avg}}$ & Backbone & L-Spar & Loc.\ Deg. & K-Neigh.\\
\midrule
com-Amazon    &  5.53 & 0.362 & 0.301 & 0.352 & 0.330\\
ca-HepTh      &  5.74 & 0.348 & 0.276 & 0.347 & 0.317\\
com-DBLP      &  6.62 & 0.302 & 0.244 & 0.301 & 0.277\\
ca-CondMat    &  8.55 & 0.234 & 0.186 & 0.234 & 0.218\\
email-Enron   & 10.73 & 0.186 & 0.159 & 0.186 & 0.179\\
wiki-Vote     & 28.51 & 0.070 & 0.067 & 0.070 & 0.069\\
email-Eu-core & 32.58 & 0.061 & 0.053 & 0.061 & 0.060\\
\bottomrule
\end{tabular}
\end{table}

\subsection{The objective}\label{sec:negative:objective}

Table~\ref{tab:negative} reports the result. Across 102 matched-retention cells,
seven sparsifiers on seven networks under a modularity optimizer, not one cell
improves modularity measured on the original graph, whether the baseline is the
compute-matched control or the best of five plain restarts. The best cell in the
entire grid loses $0.0105$.

\begin{table}[t]
\centering
\caption{Every matched-retention cell, Leiden, seven networks. $\Delta Q$ is
measured on the original graph, Eq.~\eqref{eq:transfer}, against the best of five
plain restarts. Counts are cells with a positive value. The arms differ in how
much they lose and not in whether they lose.}
\label{tab:negative}
\small
\setlength{\tabcolsep}{4pt}
\begin{tabular}{@{}lrrrr@{}}
\toprule
Arm & Cells & $\Delta Q > 0$ & Best cell & Worst cell\\
\midrule
K-Neighbor            & 18 & 0 & $-0.0105$ & $-0.667$\\
DSpar                 & 18 & 0 & $-0.0116$ & $-0.640$\\
L-Spar                &  9 & 0 & $-0.0117$ & $-0.407$\\
Backbone, random fill & 14 & 0 & $-0.0141$ & $-0.214$\\
Local Similarity      & 15 & 0 & $-0.0172$ & $-0.311$\\
Local Degree          & 14 & 0 & $-0.0288$ & $-0.278$\\
Backbone, Jaccard fill& 14 & 0 & $-0.0290$ & $-0.235$\\
\midrule
Total                 & 102 & \textbf{0} & $-0.0105$ & $-0.667$\\
\bottomrule
\end{tabular}
\end{table}

The sweep that isolates one method across its full retention range agrees. Local
similarity sparsification on the same seven networks loses between $0.014$ and
$0.404$ against a compute-matched baseline in every configuration that runs, with
the loss growing as retention falls, and its three most aggressive configurations
cannot be run at all because of the floors in Table~\ref{tab:floors}.

\subsection{The loss is not fragmentation}\label{sec:negative:fragmentation}

Sparsification breaks graphs apart, and a natural reading of the previous table is
that the optimizer is being handed a shattered graph and is penalized for it. The
backbone arms test that reading directly, because they retain a spanning structure
and so cannot increase the number of connected components.

They do what they were built to do. Under the backbone with random fill, no
detected community is a singleton in any of its fourteen cells, and the fraction
of vertices sitting in communities of fewer than three members never exceeds
$1.8\%$. The shattering arms reach $86\%$ on the same networks. The conclusion
does not move: the backbone's best cell is $-0.0141$, worse than four of the arms
that fragment freely.

The fragmentation explanation is therefore tested and rejected as sufficient. What
is left is the edges themselves. Removing half of them costs modularity on the
original graph whether or not the graph stays in one piece.

A second observation follows from the same arms and points the same way. The
backbone's two fills differ only in which edges top up the spanning structure,
random or highest Jaccard similarity. Random fill scores better in ten of the
fourteen paired cells. On these graphs, in this regime, the similarity signal
subtracts.

\subsection{The loss is not an artifact of one optimizer}\label{sec:negative:algorithms}

Everything above uses a modularity optimizer, which raises the question of whether
the objective and the algorithm are too closely coupled. Running two sparsifiers
across three further algorithms, one flow-based and one a local majority rule
alongside a second modularity optimizer, gives 63 configurations on the same seven
networks. Against the best of $N$ plain restarts, 60 are negative. The mean change
is $-0.082$ under Infomap, $-0.046$ under Louvain and $-0.106$ under label
propagation. The three positive cells are all label propagation on one network,
where the unsparsified baseline is pathological: its own best of twenty restarts
reaches modularity $0.079$ where a modularity optimizer on the same graph reaches
$0.416$. Relieving a known failure mode of one algorithm is not a benefit of
sparsification.

\subsection{Recovery, and what happens to it under control}\label{sec:negative:recovery}

Agreement with reference communities behaves differently from the objective, and
the difference is instructive rather than encouraging.

Raw agreement often rises. It rises because the sparsified partition is finer, and
both of the controls in Sections~\ref{sec:protocol:granularity}
and~\ref{sec:protocol:chance} are needed to see it. On com-DBLP, local similarity
sparsification lifts average best-match $F_1$ from $0.102$ to $0.193$ while taking
the community count from 220 to 10,947; the original graph partitioned to the same
count reaches $0.300$, and the random-partition floor at that count is $0.090$
against $0.019$ at the baseline count. Nearly all of the apparent gain is
granularity, and what is left of it is chance.

This is the pattern across the labelled networks, with one exception, which is the
subject of Section~\ref{sec:negative:exception}.

\subsection{No speed benefit either}\label{sec:negative:speed}

The pipeline is not faster. Of the 63 configurations, four preserve quality in the
sense of coming within $0.005$ of the best-of-$N$ baseline, and the fastest of
those four runs at $0.66\times$ the speed of clustering the original graph
directly. Twelve configurations exceed $1.5\times$, and every one of them loses
modularity, by $0.015$ to $0.162$. Timing only the detection step and giving the
sparsifier away for free, the median speedup is $1.21\times$, $1.63\times$ and
$1.15\times$ by algorithm, well short of the factor of two that halving the edge
count suggests, because all three algorithms return more communities on the
sparsified graph and do more work per pass.

\subsection{Hypotheses of ours that did not survive}\label{sec:negative:killed}

Four explanations for the negative result were our own, were pre-registered, and
failed.

We proposed that sparsification preserves community cores while shedding
peripheral vertices, which would make the loss benign. Four of that hypothesis's
five predictions failed: the fragments subdivide communities without mixing them,
but they cut the most embedded vertices as readily as the least, and a
resolution-matched partition of the original graph preserves cores better in
sixteen of twenty comparisons.

We proposed that the damage is concentrated in leaf and near-leaf vertices, so
that shedding them first would be harmless. It was not.

We proposed that triangle structure governs which edges a degree-based sparsifier
removes. A causal test moved the triangle count in both directions while the
score-separation quantity fell uniformly, which rules the explanation out rather
than leaving it unsupported.

Finally, we found a statistically clean gain on a large network, replicated it
exactly, and then killed it ourselves: its cluster count correlates with its
modularity at $-0.799$, the seeded run's community count falls below all twenty
plain restarts, and its mean lands on the regression line those restarts trace, so
the gain is granularity and not structure.

Reporting these is not a courtesy. Each is an explanation a reader would otherwise
have to wonder about, and each was closed by a test rather than by argument.

\subsection{The exception}\label{sec:negative:exception}

One result in this section survives every control we have.

On com-Amazon, local similarity sparsification at realized retention $0.542$
reaches average best-match $F_1$ of $0.402$ against a baseline of $0.142$. The
granularity control does not remove it: partitioning the original graph to the
same community count gives $0.319$, and deliberately over-matching by 37 per cent
more communities still gives only $0.372$. The chance floor does not remove it
either, and subtracting the floor widens the margin rather than narrowing it,
because the floor rises with the community count. A second, independent instance
appears in the same experiments from an unrelated method, Local Degree, on
com-Amazon and com-DBLP.

What removes it is a different comparison. Both non-modularity algorithms reach
higher agreement on com-Amazon without any sparsification at all, $0.465$ under
Infomap and $0.480$ under label propagation, against the $0.402$ that the
sparsified modularity pipeline attains. The effect is real, and it is the
modularity objective's default granularity being partially repaired by shattering
the graph, in a metric that rewards fine partitions. It is not a capability that
sparsification confers.

Two things about this cell are worth carrying forward. It coexists with a
modularity loss of $0.079$ in the same configuration, so the objective and the
recovery measure move in opposite directions on the same partition of the same
graph. And it is the reason the negative result of this section is stated as it
is: no benefit on the objective anywhere, and no recovery benefit that survives
comparison against a better algorithm run on the untouched graph. Section
\ref{sec:discussion} takes up the dissociation as an open phenomenon.
